\documentclass[10pt,a4paper]{amsart}
\usepackage[margin=19mm]{geometry}
\usepackage{amsmath,amssymb,mathtools}
\usepackage[scr=rsfs,cal=euler]{mathalfa}
\usepackage{xfrac}
\usepackage[unicode,hidelinks,bookmarks=false]{hyperref}
\newcommand{\ie}{i.e.\ }
\newcommand{\cf}{cf.\ }
\newcommand{\resp}{respectively\ }
\newcommand{\Subsection}{\subsection}
\providecommand{\nobreakdash}{\nobreak\hbox{-}}
\setlength{\emergencystretch}{2em}
\multlinegap0pt
\usepackage{tikz}
\usepackage{amssymb}
\newtheorem{restatable}{Restatable}
\usepackage{cleveref}
\crefname{restatable}{Restatable}{Restatables}
\newcommand{\dist}{\mathrm{dist}}
\title{On the complexity of edge subdivision to $H$-free graphs}
\author{Marta Piecyk, R. B. Sandeep}
\date{Source: arXiv:2604.24228v1 (CC BY 4.0)}
\begin{document}
\maketitle

\noindent\emph{Source and licence.} On the complexity of edge subdivision to $H$-free graphs. Version arXiv:2604.24228v1 (CC BY 4.0), distributed by arXiv under the Creative Commons Attribution 4.0 International licence, http://creativecommons.org/licenses/by/4.0/, as recorded in the arXiv licence metadata for this exact version and shown on the abstract page https://arxiv.org/abs/2604.24228v1. Full licence notice: \texttt{licenses/arxiv-2604.24228-CC-BY-4.0.txt}.\par\smallskip

\noindent\textit{Editorial scope.} Counted unit: the article's lemma OddDistanceTreeLemma (source0.tex lines 2325--2332). The article's complete original proof of it is delivered with it (source0.tex lines 2334--2491, one \texttt{proof} environment of 158 lines). No mathematics is written, repaired or re-derived: the statement and the proof are the article's own lines, in the article's own notation, and no retained line is substituted or re-flowed.  Retained uncounted as the source-local context the proof needs: lines 299--309.  Nothing was omitted from any retained span, so the package carries no omission record.  No retained line contains a citation, so the wrapper carries no bibliography and no bibliography entry.  No retained line contains a figure environment, an image-inclusion command or drawing code, so no image is delivered and the package declares no asset dependency.  Editorial formatting only, in addition: an AMS \texttt{amsart} preamble that restates the article's own declarations and macros used by the retained lines, namely the environments \texttt{restatable} on separate counters, together with the standard packages those lines need.  The retained environments print in this document as Restatable 1 in order of appearance, whereas the article prints the same items with the numbering of its own full document.  The complete native source of the article as distributed in the arXiv e-print for this exact version is bundled unchanged as \texttt{sources/199220/source0.tex}.
\begin{restatable}{theorem}{EasyTheorem}
\label{thm:easy}
Let $H$ be a graph satisfying at least one of the following properties:
\begin{itemize}
    \item $H$ has minimum degree at least $2$, and the neighborhood of every degree-$2$ vertex induces a $K_2$;
    \item the vertices of degree at least $3$ induce a graph with at least two edges.
\end{itemize}
Then \textsc{$H$-free Subdivision} is NP-complete.
Furthermore, assuming the ETH, the problem cannot be solved in time
$2^{o(k)}\cdot n^{O(1)}$.
\end{restatable}

\begin{restatable}{lemma}{OddDistanceTreeLemma}
\label{lem:tree-odd-distance-large}
Let $H$ be a tree with exactly two branching vertices $u$ and $v$.
If $\dist_H(u,v)$ is odd and at least $5$, then
\textsc{$H$-free Subdivision} is \textsf{NP}-complete.
Moreover, assuming the ETH, the problem cannot be solved in time
$2^{o(k)}\cdot n^{O(1)}$.
\end{restatable}

\begin{proof}
%Membership in NP is clear.

Let $d:=\dist_H(u,v)$. By assumption, $d$ is odd and at least $5$, and hence
\[
    d=2r+3
\]
for some integer $r\geq 1$.  Let $P$ be the unique $u$--$v$ path in $H$.
Every component of $H-V(P)$ is a path attached to either $u$ or $v$.  Let
$\lambda$ be the maximum length of such a pendant path, and let
\[
    c:=\max\{\deg_H(u)-1,\deg_H(v)-1\}.
\]
Here the term $-1$ accounts for the edge of $P$ incident with the corresponding
branching vertex.

We reduce from \textsc{$P_4$-free Edge Deletion}. %, which is \textsf{NP}-complete
%and, assuming the ETH, admits no $2^{o(k')}\cdot n^{O(1)}$-time algorithm.

Let $(G',k')$ be an instance of \textsc{$P_4$-free Edge Deletion}.  We construct
an instance $(G,K)$ of \textsc{$H$-free Subdivision}, where
\[
    K:=d k'.
\]

For each $q\in\{0,1,\ldots,r\}$, define a rooted tree $A_q$ as follows.
Start with a path of length $q$ from the root $\rho_q$ to a vertex $b_q$.
Then attach $c$ pendant paths of length $\lambda$ to $b_q$.
When $q=0$, we have $\rho_q=b_q$.

We obtain $G$ from $G'$ as follows.  For every vertex $x\in V(G')$ and every
$q\in\{0,1,\ldots,r\}$, attach $K+1$ fresh copies of $A_q$ by identifying
their roots with $x$.  All these copies are otherwise vertex-disjoint.

We claim that $(G',k')$ is a yes-instance of
\textsc{$P_4$-free Edge Deletion} if and only if $(G,K)$ is a yes-instance of
\textsc{$H$-free Subdivision}.

First suppose that there is a set $F\subseteq E(G')$ with $|F|\leq k'$ such
that $G'-F$ is $P_4$-free.  Let $J$ be the graph obtained from $G$ by subdividing
every edge of $F$ exactly $d$ times.  The number of subdivisions is
$d|F|\leq dk'=K$.

We show that $J$ is $H$-free.  Suppose, for a contradiction, that $J$ contains
an induced copy of $H$.  Let $\alpha$ and $\beta$ be the images of the two
branching vertices of $H$ in this copy.  Since vertices introduced by
subdivision have degree $2$, and since the only vertices of the attached
gadgets that can have degree at least $3$ are the original vertices of $G'$
and the vertices $b_q$, each of $\alpha$ and $\beta$ is either an original
vertex of $G'$ or a vertex $b_q$ in some attached copy of $A_q$.

For such a branching vertex $\alpha$, define its root $x_\alpha\in V(G')$ as
follows.  If $\alpha$ is an original vertex of $G'$, then set
$x_\alpha=\alpha$ and $q_\alpha=0$.  Otherwise, let $x_\alpha$ be the original
vertex to which the corresponding copy of $A_{q_\alpha}$ is attached; thus
$\dist_J(\alpha,x_\alpha)=q_\alpha$, where $0\leq q_\alpha\leq r$.
Define $x_\beta$ and $q_\beta$ analogously.

The path between $\alpha$ and $\beta$ in the induced copy of $H$ has length
exactly $d$.  Hence the part of this path between $x_\alpha$ and $x_\beta$ has
length
\[
    p=d-q_\alpha-q_\beta .
\]
Since $q_\alpha,q_\beta\leq r$, we have
\[
    p\geq d-2r=3,
\]
and clearly $p\leq d$.

No edge of $F$ can occur on this path between $x_\alpha$ and $x_\beta$, because
every edge of $F$ was replaced by a path of length $d+1$, while the whole path
between $x_\alpha$ and $x_\beta$ has length at most $d$.  Therefore this path
corresponds to an induced path of length at least $3$ in $G'-F$.  This is
impossible because $G'-F$ is $P_4$-free; indeed, any induced path of length at
least $3$ contains an induced $P_4$.  Hence $J$ is $H$-free.

Conversely, suppose that $G$ can be transformed into an $H$-free graph $J$ by
at most $K$ subdivisions.  For each edge $e\in E(G')$, let $s(e)$ be the number
of times $e$ is subdivided in this solution.  Define
\[
    F:=\{e\in E(G') : s(e)\geq d\}.
\]
Since the total number of subdivisions is at most $K=dk'$, we have
\[
    |F|\leq k'.
\]

We claim that $G'-F$ is $P_4$-free.  Suppose not, and let $wxyz$ be an induced
$P_4$ in $G'-F$.  Let
\[
    p_1=s(wx)+1,\qquad p_2=s(xy)+1,\qquad p_3=s(yz)+1.
\]
Since none of $wx,xy,yz$ belongs to $F$, we have
\[
    1\leq p_i\leq d
    \quad\text{for each }i\in\{1,2,3\}.
\]

We first observe that among the consecutive sums
\[
    p_i,\quad p_i+p_{i+1},\quad p_1+p_2+p_3,
\]
there is one whose value lies between $3$ and $d$.  Indeed, if some
$p_i\geq 3$, then $p_i\in[3,d]$.  Otherwise all $p_i\leq 2$.  If
$p_1=p_2=p_3=1$, then $p_1+p_2+p_3=3$.  If some $p_i=2$, then adding an
adjacent $p_j$ gives a consecutive sum between $3$ and $4$, which is at most
$d$ because $d\geq 5$.

Thus, in the graph $J$, the subdivided version of the induced path $wxyz$
contains an induced path $Q$ of length $p$ between two original vertices
$a,b\in\{w,x,y,z\}$, where
\[
    3\leq p\leq d.
\]
The path $Q$ is induced: the path $wxyz$ is induced in $G'-F$, and any edge
between nonconsecutive vertices of this path that belongs to $F$ has been
replaced by a path of length at least $d+1$, so it gives no chord in $J$.

Since $3\leq p\leq d=2r+3$, we have
\[
    0\leq d-p\leq 2r.
\]
Therefore, we can choose integers $q_a,q_b\in\{0,1,\ldots,r\}$ such that
\[
    q_a+q_b=d-p.
\]

For each fixed original vertex and each fixed value of $q$, there are $K+1$
attached copies of $A_q$, while the solution uses at most $K$ subdivisions in
total.  Hence there is an untouched copy of $A_{q_a}$ attached to $a$ and an
untouched copy of $A_{q_b}$ attached to $b$.

Let $\alpha$ and $\beta$ be the vertices $b_{q_a}$ and $b_{q_b}$ in these two
untouched copies.  Then the path from $\alpha$ to $\beta$ obtained by going
from $\alpha$ to $a$, then along $Q$ from $a$ to $b$, and finally from $b$ to
$\beta$ has length
\[
    q_a+p+q_b=d.
\]
Moreover, the two untouched gadgets contain enough pendant paths of length at
least those required at the two branching vertices of $H$: the vertex
$\alpha$ has $c$ pendant paths of length $\lambda$, and so does $\beta$.
By choosing the appropriate number of these pendant paths and taking suitable
initial subpaths, we obtain an induced copy of $H$ in $J$.  This contradicts
the assumption that $J$ is $H$-free.

Therefore $G'-F$ is $P_4$-free.  Since $|F|\leq k'$, the instance $(G',k')$ is
a yes-instance of \textsc{$P_4$-free Edge Deletion}.

This completes the correctness of the reduction. Now, the statements follow from \cref{thm:easy}.
%The reduction is polynomial-time, and the new parameter is
%$K=dk'$, where $d$ depends only on the fixed graph $H$.  Hence an algorithm for
%\textsc{$H$-free Subdivision} running in time
%$2^{o(k)}\cdot n^{O(1)}$ would imply an algorithm for
%\textsc{$P_4$-free Edge Deletion} running in time
%$2^{o(k')}\cdot n^{O(1)}$, contradicting the ETH.
\end{proof}

\end{document}
